\chapter{Родитель происхождения
\texorpdfstring{$M_4$}{M4} cross-генератора}
\label{ch:v10-hopf-reservoir-m4-cross-parent}

\section{Минимальный динамический кандидат}

На допущенном общем носителе cross-блок задаётся матрицей
$V\in M_2(\mathbb C)$. Минимальный положительный функционал, одновременно
требующий частичной изометрии и сохранения ориентации, имеет вид
\begin{equation}
 S_V=\frac14\Tr(V^*V-I)^2
 +\frac14\Tr[(ZV-VZ)^*(ZV-VZ)],
 \qquad Z=\operatorname{diag}(-1,1).
 \label{eq:v10-m4-cross-parent}
\end{equation}
На вещественном срезе в точке $V=I_2$ градиент равен нулю, а гессиан в
координатах $(v_{11},v_{12},v_{21},v_{22})$ равен
\begin{equation}
 H_V=\begin{pmatrix}
 2&0&0&0\\
 0&3&1&0\\
 0&1&3&0\\
 0&0&0&2
 \end{pmatrix},qquad
 Spec H_V=\{2,2,2,4\},\qquad \det H_V=32.
\end{equation}
Следовательно, условный родитель действительно может породить устойчивый
ненулевой cross-блок.

\section{Норма выбрана, фаза не выбрана}

Функционал имеет четыре вещественных нулевых минимума
\begin{equation}
 V=\operatorname{diag}(\pm1,\pm1).
\end{equation}
В комплексном секторе вырождение больше:
\begin{equation}
 V(\phi_h,\phi_c)=
 \operatorname{diag}(e^{i\phi_h},e^{i\phi_c}),\qquad
 S_V[V(\phi_h,\phi_c)]=0.
\end{equation}
Таким образом, ориентационный член запрещает смешение горячего и холодного
уровней, но не выбирает две диагональные фазы.

Кроме того, член $V^*V-I$ уже содержит выбранную ненулевую норму
конденсата. В общей ландауовской записи она определяется отношением
квадратичного и квартального коэффициентов; текущий корпус не выводит ни
их знак, ни отношение.

\section{Унаследованный родитель}

До допуска $M_4$ переменная $V$ отсутствует. Если формально добавить её с
нулевыми коэффициентами, унаследованные cross-гессиан и линейный источник
равны
\begin{equation}
 H_V^{\rm inherited}=0_{4\times4},\qquad j_V^{\rm inherited}=0,
 \qquad \rank H_V^{\rm inherited}=0.
\end{equation}
Независимая симметрия двух секторов не допускает ненулевого линейного
бифундаментального источника. Поэтому~\eqref{eq:v10-m4-cross-parent}
является строгим условным завершением, но пока не следствием предыдущего
родителя.

\begin{theorem}[Условная cross-конденсация]
Квартальный partial-isometry-потенциал с ориентационным штрафом имеет
устойчивые ненулевые минимумы и способен условно породить генератор,
расширяющий $M_2\oplus M_2$ до $M_4$.
\end{theorem}

\begin{nogocriterion}[Конденсат без источника и фазового селектора]
Положительность гессиана в заранее нормированной точке $V=I_2$ не является
происхождением cross-генератора. Требуются независимо выведенные переменная
$V$, знак квадратичного члена, масштаб конденсата и член, снимающий
$U(1)^2$-вырождение без загрузки целевого интертвинера.
\end{nogocriterion}

\begin{protocol}\begin{enumerate}
\item условная архитектура: \textbf{$10/10$};
\item условное алгебраическое происхождение: \textbf{$8/8$};
\item унаследованный cross-гессиан: \textbf{ранг $0/4$};
\item происхождение коэффициентов: \textbf{$0/3$};
\item строгое физическое происхождение: \textbf{$0/4$};
\item следующий гейт: \textbf{аудит кандидатов бифундаментального конденсата}.
\end{enumerate}\end{protocol}

Машинный сертификат сохранён в
\path{s2t_v10_cell_birth_four_volume_induced_newton_breathing_anomaly_hopf_reservoir_intertwiner_m4_cross_generator_parent_origin_gate_results.json}.